% ==========================================
% SEÇÃO 4: MATEMÁTICA E PROCESSAMENTO
% ==========================================
\section[Processamento]{Matemática e Processamento}

\begin{frame}{Transformações de Espaço e o Histograma}
    \begin{itemize}
        \item \textbf{Matrizes de Transformação:}
        \begin{itemize}
            \item A conversão entre espaços de cor não é arbitrária; ela ocorre via multiplicação de matrizes por vetores de pixels $[R, G, B]^T$.
            \item \textbf{CIE XYZ:} Espaço matemático universal usado como "ponte" de tradução entre hardwares diferentes.
        \end{itemize}
        
        \vspace{0.2cm}
        
        \item \textbf{O Histograma como Array de Dados:}
        \begin{itemize}
            \item É um vetor de contagem estatística discreta. Para imagens de 8 bits, cria-se um array de 256 posições (\textit{bins}).
            \item Varre-se a matriz de pixels incrementando o índice correspondente à intensidade encontrada.
        \end{itemize}
        
        \vspace{0.2cm}
        
        \item \textbf{Engenharia de Algoritmos (CBIR):}
        \begin{itemize}
            \item Sistemas de busca reversa de imagens comparam a distância euclidiana ou a interseção entre vetores de histogramas para medir similaridade cromática.
        \end{itemize}
    \end{itemize}

    \note{
        \begin{columns}[T]
            \begin{column}{0.49\textwidth}
                \scriptsize
                \textbf{Pipeline de Álgebra Linear:}
                \begin{itemize}
                    \setlength{\itemsep}{1pt}
                    \item \textbf{Pipeline de GPU:} A conversão RGB para XYZ utiliza uma matriz $3 \times 3$ fixa com coeficientes baseados no observador padrão. Como essa operação é uma transformação linear pura e independente para cada pixel, ela é paralelizada nativamente em Shaders dentro da GPU, permitindo processamento em tempo real de fluxos de vídeo em alta resolução.
                \end{itemize}
            \end{column}

            \begin{column}{0.49\textwidth}    
                \scriptsize
                \textbf{Manipulação Estatística:}
                \begin{itemize}
                    \setlength{\itemsep}{1pt}
                    \item \textbf{Equalização de Histograma:} Algoritmos de melhoria de contraste usam a Função de Distribuição Acumulada (CDF) do histograma para redistribuir as frequências de pixels. Áreas muito concentradas (subexpostas ou superexpostas) têm seus valores espalhados linearmente por todo o array de 0 a 255, maximizando a entropia e o contraste da imagem de forma automatizada.
                \end{itemize}
            \end{column}
        \end{columns}   
    }
\end{frame}